\section*{第 \ref{ch:b3:endocrinology} 章——内分泌与稳态}

\begin{solution}{exo:b3:endocrinology:1}
\omterm{def:b3:endocrinology:glucose}{胰岛素}：肽类，表面受体（一种酪氨酸激酶），起效时间为
数分钟，半衰期约 \qty{5}{min}。皮质醇：类固醇，\omterm{def:b3:endocrinology:hormone}{核受体}，
起效时间为数小时，半衰期 \qty{80}{min}（大部分与载体蛋白结合）。
肾上腺素：胺类，表面 G 蛋白偶联受体，起效时间为数秒，
半衰期约 \qty{2}{min}。甲状腺素：碘化氨基酸，\omterm{def:b3:endocrinology:hormone}{核受体}，
起效时间为数天，半衰期 \qty{7}{d}（与载体蛋白结合）。
\omterm{prop:b3:renal-osmoregulation:countercurrent}{抗利尿激素}：肽类，表面 G 蛋白偶联受体（经由 \omterm{prop:b3:learning-memory:aplysia}{cAMP}，在
\omterm{def:b3:renal-osmoregulation:nephron}{集合管}中起作用），起效时间为数分钟，半衰期约 \qty{15}{min}。
\end{solution}

\begin{solution}{exo:b3:endocrinology:2}
\omterm{def:b3:endocrinology:axes}{下丘脑}（TRH）$\to$ \omterm{def:b3:endocrinology:axes}{垂体}（TSH）$\to$ 甲状腺（T$_{4}$），其中
T$_{4}$ 同时抑制上两级。甲状腺被破坏：T$_{4}$ 低，
TSH 高。分泌 TSH 的\omterm{def:b3:cancer-biology:hallmarks}{肿瘤}：TSH 高，T$_{4}$ 也高（反馈
无法抑制该\omterm{def:b3:cancer-biology:hallmarks}{肿瘤}）。过量服用甲状腺素片：
T$_{4}$ 高，TSH 被抑制。缺碘：T$_{4}$ 低或
偏低，TSH 高，腺体在其刺激下增大为\omterm{prop:b3:endocrinology:thyroid}{甲状腺肿}。
\end{solution}

\begin{solution}{exo:b3:endocrinology:3}
大脑在数分钟内因血糖过低而死亡，血糖过高则
要经过数年才造成伤害；而且演化过程中饥荒远比饱食常见。
因此，血糖下降时的防御是冗余的（\omterm{def:b3:endocrinology:glucose}{胰高血糖素}、肾上腺素、
皮质醇、生长\omterm{def:b3:endocrinology:hormone}{激素}），而血糖上升时的防御是单一的。
由此可作预测：\omterm{def:b3:endocrinology:glucose}{胰岛素}过多会急性致死（低血糖昏迷），
\omterm{def:b3:endocrinology:glucose}{胰岛素}过少则是慢性疾病——这正是临床所见。
\end{solution}

\begin{solution}{exo:b3:endocrinology:4}
\omterm{def:b3:endocrinology:clock}{授时因子}是使生物钟同步的环境线索：光照、
进食时间、温度、运动和社会作息。主生物钟
利用光照，光信号经\omterm{def:b3:sensory-systems:retina}{视网膜}中的黑视蛋白\omterm{def:b3:sensory-systems:retina}{神经节细胞}传入。穿越
八个时区后，生物钟每天只调整约一小时，因此一周之内，
睡眠、皮质醇、体温和食欲仍按原来的时间运行，
而由进食重新设定的外周生物钟则各自按自身节奏漂移：
这就是时差反应；向东旅行时更严重，因为把生物钟提前
比推迟更困难。
\end{solution}

\begin{solution}{exo:b3:endocrinology:5}
占据 $250$ 个受体时达到半最大响应：$H = K_{d}\times
250/(\num{20000} - 250) = \qty{1}{nM}\times 0.0127 = \qty{12.7}{pM}$，
比 $K_{d}$ 低八十倍。若只有 \num{2000} 个受体：$250/1750
\times \qty{1}{nM} = \qty{143}{pM}$，敏感性降低为原来的十一分之一。
长期高浓度\omterm{def:b3:endocrinology:glucose}{胰岛素}会使其受体\omterm{prop:b3:endocrinology:dose}{下调}，因此要达到同样的效应需要更多\omterm{def:b3:endocrinology:glucose}{胰岛素}，
这是\omterm{def:b3:endocrinology:glucose}{胰岛素}抵抗的一个组成部分，它
会自我强化。
\end{solution}

\begin{solution}{exo:b3:endocrinology:6}
$T_{4} = 13$：$1.5\,\mathrm{e}^{0.92} = \qty{3.8}{mU/L}$；$11$：
$1.5\,\mathrm{e}^{1.84} = \qty{9.4}{mU/L}$；$9$：$1.5\,\mathrm{e}^{2.76}
= \qty{24}{mU/L}$；$25$：$1.5\,\mathrm{e}^{-4.6} = \qty{0.015}{mU/L}$。在
\qty{11}{pmol/L} 时，T$_{4}$ 仍在 \qtyrange{10}{20}{} 范围内，
而 TSH 超过其上限的两倍（亚临床\omterm{prop:b3:endocrinology:thyroid}{甲状腺功能减退}）；
在 $13$ 时 TSH 刚好在正常范围内；在 $9$ 时两者均异常；在 $25$ 时
两者均朝相反方向异常。
\end{solution}

\begin{solution}{exo:b3:endocrinology:7}
$L = s\beta/(a\gamma) = 0.2\times 0.05/(0.02\times 0.1) = 5$。稳态过量
$g^{*} = (u/a)/(1 + L) = (2/0.02)/6 = \qty{16.7}{mg/dL}$；无反馈时
$u/a = \qty{100}{mg/dL}$。\omterm{def:b3:endocrinology:glucose}{胰岛素}过量 $i^{*} = \beta g^{*}/
\gamma = 0.05\times 16.7/0.1 = \qty{8.3}{mU/L}$。
\end{solution}

\begin{solution}{exo:b3:endocrinology:8}
判别式 $(a - \gamma)^{2} - 4s\beta = 0.0064 - 0.04 < 0$：
为振荡型。$\omega = \sqrt{a\gamma + s\beta - (a+\gamma)^{2}/4} =
\sqrt{0.002 + 0.01 - 0.0036} = \qty{0.092}{min^{-1}}$，周期 $2\pi/
\omega = \qty{68}{min}$；振幅衰减到原来的 $\mathrm{e}$ 分之一所需时间为 $2/(a +
\gamma) = \qty{16.7}{min}$。将 $s$ 减半：$s\beta = 0.005$，判别式
$0.0064 - 0.02 < 0$，仍为振荡型，$\omega = \sqrt{0.007 - 0.0036}
= \qty{0.058}{min^{-1}}$，周期 \qty{108}{min}；衰减时间不变；
$L = 2.5$。较弱的环路回复得更慢，
并保持较大的稳态误差。
\end{solution}

\begin{solution}{exo:b3:endocrinology:9}
一年来，泼尼松龙抑制了 CRH 和 ACTH 的分泌，
未受刺激的肾上腺皮质随之萎缩。停药后，体内已无外源类固醇，
\omterm{def:b3:endocrinology:axes}{下丘脑}和\omterm{def:b3:endocrinology:axes}{垂体}需要数周才能恢复，
而萎缩的肾上腺即使接收到 ACTH 也无法作出反应。感染时需要的皮质醇是正常量的十倍，
若无足够皮质醇，血管扩张，血压和血糖下降：即肾上腺危象。
检测结果为：ACTH 低（病变位于中枢），皮质醇低，
注射 ACTH 后皮质醇升高很少（腺体已萎缩），
这与 ACTH 升高的阿狄森病不同。
因此需要缓慢减量。
\end{solution}

\begin{solution}{exo:b3:endocrinology:10}
单变量情形：$\dot x = f(x)$ 且 $f' < 0$，
只有一个不动点 $x^{*}$；当 $x > x^{*}$ 时 $\dot x < 0$，
$x$ 单调减小并趋向 $x^{*}$，且不会越过它（解的唯一性），
下方情形对称，因此不会出现振荡。
双变量情形：当效应器通过一个具有自身时间尺度的第二变量起作用时，
矩阵可以有复特征值，如定理所述。
瞬时的抑制会稳定在某一水平上；转录与核内抑制之间的延迟
（包括翻译、二聚化、磷酸化和入核）使蛋白质在基因被关闭之前超调，
在基因重新开启之前欠调，
周期大约为延迟与蛋白质降解时间之和的两倍。
提高 PER 的降解速率会缩短周期：
人 PER2 的一个降低稳定性的突变使生物钟周期约为 20 小时，
并使一个家族在晚上 7 点入睡。
\end{solution}

\begin{solution}{exo:b3:endocrinology:11}
指数为 $k(\text{Ca} - 1.2)$：在 $1.10$ 时为 $\mathrm{e}^{-2}$，PTH $= 1/
(1 + 0.135) = 0.88$；$1.15$：$0.73$；$1.20$：$0.50$；$1.25$：$1/(1 +
2.72) = 0.27$；$1.30$：$1/(1 + 7.39) = 0.12$。
在设定点处的斜率为 $-k/4 = -5$（每 mmol/L），
即每 \qty{0.01}{mmol/L} 变化 \qty{5}{\%} 的最大响应。
钙必须维持在几个百分点之内，其效应器作用于巨大的缓冲库（骨骼）且不发生超调，
因此陡峭的响应既安全又必要。
若血糖的环路如此陡峭，并通过具有自身清除时间的\omterm{def:b3:endocrinology:glucose}{胰岛素}起作用，
它将具有很大的\omterm{thm:b3:endocrinology:loop}{环路增益}，并在每次进餐后振荡至深度低血糖。
\end{solution}

\begin{solution}{exo:b3:endocrinology:12}
糖基化不可逆且缓慢，因此年龄为 $\tau$ 的细胞
中被糖基化的比例为 $kG\tau$；对年龄均匀分布于
\numrange{0}{120}~天之间的细胞取平均，\omterm{met:b3:endocrinology:ogtt}{HbA1c} $= kG\times\qty{60}{d}$，与平均血糖成正比。
校准：$5.5\times 60k = 0.05$，得 $k =
\qty{1.5e-4}{}$ 每（mmol/L）·天；在 \qty{10}{mmol/L} 时：\qty{9.1}{\%}。
若细胞寿命为 60 天，则平均年龄为 30 天，\omterm{met:b3:endocrinology:ogtt}{HbA1c}
在血糖相同时减半：血糖为 \qty{10}{mmol/L} 的患者，其读数为
\qty{4.5}{\%}，属正常。
\end{solution}

\begin{solution}{pb:b3:endocrinology:1}
\textbf{1.} \qty{90}{mg/dL} $= \qty{0.9}{g/L} = 0.9/180 = \qty{5.0}{mmol/L}$;
$0.9\times 15 = \qty{13.5}{g}$ 存在于分布容积中。
\textbf{2.} $75/15 = \qty{5}{g/L}$：升高 \qty{500}{mg/dL}，
\qty{28}{mmol/L}。实际峰值远低于此，因为吸收
分布在一小时内完成，清除随葡萄糖到达即刻进行，且
肝脏在首过时摄取了很大一部分。
\textbf{3.} $\qty{75}{g}/\qty{60}{min} = \qty{1.25}{g/min}$；在
\qty{15}{L} 的分布容积内：\qty{8.3}{mg/dL/min}，是 $u_{0}$ 的四倍。
\textbf{4.} 稳态过量 $u/a = 8.3/0.02 = \qty{417}{mg/dL}$，时间
常数 $1/a = \qty{50}{min}$；一小时后 $417(1 -
\mathrm{e}^{-1.2}) = \qty{291}{mg/dL}$，高于空腹值，约为
\qty{380}{mg/dL}，即 \qty{21}{mmol/L}。
\textbf{5.} $L = 0.2\times 0.05/(0.02\times 0.1) = 5$；稳态过量
$417/6 = \qty{69}{mg/dL}$，为无\omterm{def:b3:endocrinology:glucose}{胰岛素}时的六分之一。
\textbf{6.} 脑每天消耗 \qty{120}{g} 葡萄糖，自身不储存，
几乎无法利用其他能源：血糖过低，数分钟内即
可致意识混乱和昏迷。血糖过高则使蛋白质糖化，经数年，\omterm{def:b3:dna-repair:lesions}{损伤}
\omterm{def:b3:sensory-systems:retina}{视网膜}、肾脏和神经的小血管，以及大动脉。
\textbf{7.} $\lambda^{2} + (a + \gamma)\lambda + (a\gamma + s\beta) = 0$:
迹 $-0.12$，行列式 $0.002 + 0.01 = 0.012$，判别式
$0.0144 - 0.048 = -0.0336$.
\textbf{8.} $\omega = \sqrt{0.0336}/2 = \qty{0.092}{min^{-1}}$；周期
\qty{68}{min}；衰减时间 $2/0.12 = \qty{16.7}{min}$。
\textbf{9.} $\dot g(0) = 100(-0.06 + c\,\omega) = -2$，故 $c\,\omega =
0.04$, $c = 0.44$.
\textbf{10.} $g = 0$ 当 $\tan\omega t = -1/c = -2.29$，$\omega t =
\pi - 1.16 = 1.98$，$t = \qty{21.6}{min}$。第一个极小值处，约在
\qty{32}{min}: $100\,\mathrm{e}^{-1.94}(\cos 2.97 + 0.44\sin 2.97) =
14.4\times(-0.91) \approx -\qty{13}{mg/dL}$，血糖为 \qty{77}{mg/dL}。
\textbf{11.} 当 $t = 0$ 时，$\dot i = 100\beta > 0$：\omterm{def:b3:endocrinology:glucose}{胰岛素}在
$g > 0$ 时上升，并在 $\beta g = \gamma i$ 时达到峰值，这要求 $g$ 仍
为正——因此峰值早于血糖降至空腹水平，且在
血糖开始下降之后（它自 $t = 0$ 起即下降）。模型中峰值
约在 \qty{12}{min}。
\textbf{12.} 若输入分布在一小时内，峰值出现在
吸收与清除达到平衡之时，峰值更低、更晚（30--60 分钟）；
回落需等待吸收结束，因此需两小时；\omterm{def:b3:endocrinology:glucose}{胰岛素}
超调因输入渐进而较小，故血糖低谷较轻。
\textbf{13.} $L = 2.5$。健康者 $g^{*} = (2/0.02)/6 = \qty{16.7}{mg/dL}$；
\omterm{def:b3:endocrinology:glucose}{胰岛素}抵抗者 $100/3.5 = \qty{28.6}{mg/dL}$：升高约
\qty{12}{mg/dL}（\qty{0.7}{mmol/L}），空腹血糖约为
\qty{102}{mg/dL}，即 \qty{5.7}{mmol/L}——空腹血糖受损
的阈值。
\textbf{14.} $i^{*} = \beta g^{*}/\gamma$：健康者 $8.3$，\omterm{def:b3:endocrinology:glucose}{胰岛素}抵抗者
\qty{14.3}{mU/L}，比值 $1.7$：代偿性\omterm{def:b3:endocrinology:glucose}{胰岛素}抵抗，
伴高\omterm{def:b3:endocrinology:glucose}{胰岛素}血症，血糖接近正常。
\textbf{15.} $L = 0.1\times 0.01/0.002 = 0.5$；$g^{*} = 100/1.5 =
\qty{66.7}{mg/dL}$，高于健康状态 \qty{50}{mg/dL}：
\qty{2.8}{mmol/L}，空腹血糖 \qty{7.8}{mmol/L}（\qty{140}{mg/dL}），
高于\omterm{def:b3:endocrinology:glucose}{糖尿病}阈值。
\textbf{16.} 判别式 $(a - \gamma)^{2} - 4s\beta = 0.0064 - 0.004 =
0.0024 > 0$：单调回归。$\lambda = -0.06 \pm 0.0245$：$-0.0355$
和 $-0.0845$ 每分钟；最慢时间常数为 \qty{28}{min}，
与健康衰减的 \qty{16.7}{min} 相比：失效环路回归
更慢，且从不低于基线。
\textbf{17.} 血糖与\omterm{def:b3:endocrinology:glucose}{胰岛素}均高：伴部分代偿的\omterm{def:b3:endocrinology:glucose}{胰岛素}抵抗
——属 2 型；两小时值 \qty{13}{mmol/L} 超过
\qty{11.1}{}，故为\omterm{def:b3:endocrinology:glucose}{糖尿病}，而非仅为糖耐量受损。
\textbf{18.} 无 C 肽即无内源性\omterm{def:b3:endocrinology:glucose}{胰岛素}：属 1 型。缺乏
\omterm{def:b3:endocrinology:glucose}{胰岛素}时，肌肉与脂肪无法摄取葡萄糖，脂肪被分解为
脂肪酸和酮体，肌肉蛋白被分解以供糖异生，
葡萄糖随尿液流失并带走其热量：
患者身处富足之中却饥饿。
\textbf{19.} $5\times 9/5.5 = \qty{8.2}{\%}$.
\textbf{20.} $1.5\,\mathrm{e}^{0.46\times 3} = 1.5\times 3.97 =
\qty{6.0}{mU/L}$。T$_{4}$ 在正常范围内，TSH 高于上限：亚临床
\omterm{prop:b3:endocrinology:thyroid}{甲状腺功能减退}，腺体功能衰退，\omterm{def:b3:endocrinology:axes}{垂体}正在代偿。
\textbf{21.} $1.5\,\mathrm{e}^{0.46\times 8} = 1.5\times 39.6 =
\qty{59}{mU/L}$。T$_{4}$ 偏低而 TSH 仅为 $0.3$，表明
\omterm{def:b3:endocrinology:axes}{垂体}未作出应答：病变位于中枢（\omterm{def:b3:endocrinology:axes}{垂体}或
\omterm{def:b3:endocrinology:axes}{下丘脑}），而不在甲状腺。
\textbf{22.} 一个半衰期：水平降至 \qty{50}{\%}，而
症状在数周内悄然出现。皮质醇的半衰期为
\qty{80}{min}，漏服后数小时内即消失，当日若遇应激
则无防备：漏服一周的甲状腺素
令人不适，而漏服一天的皮质醇则可致命。
\textbf{23.} T$_{4}$ 偏高（\omterm{def:b3:adaptive-immunity:antibody}{抗体}驱动腺体，不受
反馈影响）；高 T$_{4}$ 使 TSH 受抑制；对数线性关系
使其降至不可检测水平——当 $T_{4} = 30$ 时，$1.5\,\mathrm{e}^{-6.9}
= \qty{0.0015}{mU/L}$——因此 TSH 不可检测是首个征兆。
\textbf{24.} \omterm{def:b3:endocrinology:hormone}{激素}水平反映其指令与所属腺体之间的
平衡：TSH 高而 T$_{4}$ 低为腺体失效，TSH 高而 T$_{4}$ 高
则为\omterm{def:b3:endocrinology:axes}{垂体}失控；高\omterm{def:b3:endocrinology:glucose}{胰岛素}伴高血糖提示
\omterm{def:b3:endocrinology:glucose}{胰岛素}抵抗，低\omterm{def:b3:endocrinology:glucose}{胰岛素}伴高血糖则为胰岛被毁。
单独看，T$_{4}$ 为 \qty{7}{pmol/L} 或\omterm{def:b3:endocrinology:glucose}{胰岛素}为
\qty{30}{mU/L}，并不能说明故障位于何处；
与其调控\omterm{def:b3:endocrinology:hormone}{激素}一并看，才能判断。
\textbf{25.} \omterm{thm:b3:endocrinology:loop}{环路增益} $L = 5$；周期 \qty{68}{min}，衰减时间
\qty{17}{min}；失效环路的空腹血糖 \qty{7.8}{mmol/L}。
\end{solution}
